\pdfoutput=1
\documentclass[11pt]{amsart}
\usepackage{amsmath,amssymb,amsthm}
\usepackage[margin=1.2in]{geometry}

\usepackage[colorlinks=true,linkcolor=blue,citecolor=blue,urlcolor=blue]{hyperref}
\hypersetup{pdfauthor={Bernd Johannes Wuebben},pdftitle={Local equations for deformations of Calabi-Yau threefolds with del Pezzo cone points}}
\setlength{\emergencystretch}{2em}
\newtheorem{theorem}{Theorem}[section]
\newtheorem{proposition}[theorem]{Proposition}
\newtheorem{lemma}[theorem]{Lemma}
\newtheorem{corollary}[theorem]{Corollary}
\theoremstyle{definition}
\newtheorem{example}[theorem]{Example}
\theoremstyle{remark}
\newtheorem{remark}[theorem]{Remark}
\newcommand{\C}{\mathbb C}
\newcommand{\Q}{\mathbb Q}
\newcommand{\Z}{\mathbb Z}
\newcommand{\OO}{\mathcal O}
\newcommand{\PP}{\mathbb P}
\DeclareMathOperator{\Def}{Def}
\DeclareMathOperator{\Ext}{Ext}
\DeclareMathOperator{\Spec}{Spec}
\DeclareMathOperator{\im}{im}
\DeclareMathOperator{\rk}{rank}
\title[Local equations for Calabi--Yau deformations]{Local equations for deformations of Calabi--Yau threefolds with del Pezzo cone points}
\author{Bernd Johannes Wuebben}
\date{November 26, 2026}
\subjclass[2020]{14J32 (primary); 14B07, 14D15, 14E15 (secondary)}
\keywords{Calabi--Yau threefold, deformation space, del Pezzo surface, partial smoothing, vanishing period}
\begin{document}
\begin{abstract}
We determine the analytic deformation space of a projective Calabi--Yau
threefold with ordinary double points and anticanonical cones over smooth
del Pezzo surfaces of degrees five and six, admitting a smooth projective
crepant resolution small at the nodes and given by the vertex blow-up at
the cones. Up to a smooth factor, its semiuniversal base is the
scheme-theoretic intersection of the product of the marked local bases with
the space of homological relations on the resolution. Link periods identify
the strata on which prescribed singularities persist. We construct two
projective examples with the same smooth reduced germ and the same
tangent dimension. The nilradical has square zero in one and nonzero
square in the other. The proof establishes compatibility of link periods with
Artin base change and detects global obstructions to lifting by
localization and pairing with Cartier divisor classes.
\end{abstract}
\maketitle
\section{Introduction}
Friedman's criterion for a nodal Calabi--Yau threefold asks for a
homological relation among exceptional curves with nonzero coefficient
at every node \cite{Fri86}. In a projective conifold transition the
individual nodal discriminants are linear hyperplanes in vanishing-period
coordinates \cite[Lemma~3.13 and Proposition~3.16]{LLW}. A cone over a
degree-six del Pezzo surface presents a further issue. Its local
deformation space has two components, a line and a plane. A global
first-order deformation can have nonzero projections to both at one
vertex. Such a direction lies on neither component and can occur even
when all deformations over reduced bases are locally trivial.

Let $X$ be a connected normal projective complex threefold with
$\omega_X\simeq\OO_X$ and $H^1(X,\OO_X)=0$. Its singularities are
ordinary double points and germs analytically isomorphic to the
anticanonical cones over smooth del Pezzo surfaces of degrees five and
six. We assume a smooth projective crepant resolution $\pi:Y\to X$,
small at the nodes and given by the blow-up of each cone vertex. Fix the
analytic identifications of the singularities, the markings of their
exceptional curves and the corresponding choices of small resolution
in the local modifications which replace a degree-six vertex by two
or three nodes (Section~\ref{sec:periods}).

For each singular point $p$ let $L_p$ be its oriented link and put
$V_p=H_2(L_p;\C)$. The map to the local resolution identifies $V_p$
with the signed exceptional-curve line at a node and, at a vertex, with
the perpendicular to the canonical class of the exceptional surface $E_p$
in $H_2(E_p;\C)$. A reversal of a nodal marking changes the sign of that
coordinate. Set
\begin{equation}\label{eq:kernel}
 V=\bigoplus_p V_p,\qquad
 K=\ker\left(V\longrightarrow H_2(Y;\C)\right),\qquad
 h=h^1(X,\Theta_X).
\end{equation}
The arrow is induced by the inclusions of the links and adds the marked
classes. At a node or a degree-five vertex put $G_p=V_p$. At a degree-six
vertex fix the cyclically ordered six boundary curves $B_0,\ldots,B_5$
of $E_p$. In a compatible marking put
\[
 V_{p,L}=\C(B_0-B_3),\qquad
 V_{p,P}=\Big\{\sum_{i=0,2,4}c_i B_i:\ \sum c_i=0\Big\}.
\]
These complementary summands have dimensions one and two. Write
$G_p=V_{p,L}\cup V_{p,P}$ with its reduced structure, and $G=\prod_pG_p$.
Thus $G$ is a fixed closed subscheme of $V$. The local marked
discriminants in these coordinates are described below.

\begin{theorem}\label{thm:main}
With these hypotheses and markings there is an isomorphism of complex
analytic germs
\[
 (\Def(X),0)\ \simeq\
 \big(\C^h\times (K\times_V G),0\big).
\]
Its projection to $K$ is the tuple of marked link periods of one relative
volume form, with each cohomology class identified with its Poincar\'e
dual. It identifies the inverse image of each product of local components
and of its marked discriminant hyperplanes with the corresponding
scheme-theoretic linear sections.
\end{theorem}

In coordinates $(s_p,u_p,v_p)$ in which the local line and plane are
respectively $u_p=v_p=0$ and $s_p=0$, the ideal on $K$ is generated by
the restrictions of $s_pu_p$ and $s_pv_p$ at every degree-six vertex.
The smooth coordinates along $\C^h$ are auxiliary. Replacing the global
relative generator by a base unit multiplies all its period coordinates
by that unit. The assertion concerns the marked semiuniversal space.

The two degree-six branches and their scheme structure are due to
Altmann \cite{Alt}. For the local models considered here, the image of
global tangent localization is computed in \cite[Theorem~2.5]{W5};
the existence of full smoothings is characterized in
\cite[Theorem~4.7]{W5}. Theorem~5.3 of the same paper determines a
reduced four-component unrestricted germ in a two-cone example.
For nodes and degree-six cones,
\cite[Theorems~2.2 and~3.1]{Partial} proves smoothness and determines
all marked vanishing-period divisors after a component is chosen at
each cone. In Theorem~\ref{thm:main} the same coordinates apply across
all choices and determine the ideal of the unrestricted germ. In
particular the local equations retain their meaning for an Artin family
which belongs to neither branch at a given cone.

Namikawa's stratification for terminal Calabi--Yau threefolds uses
projective small modifications \cite[Proposition~2.3 and
Theorems~2.5,~2.7]{Nam}. A crepant exceptional surface places the cone
points here outside the terminal case. Gross proves a local-to-global
comparison for canonical threefolds and gives nonreduced examples
\cite[Theorems~1.16,~2.2 and Example~2.19]{Gross}. Generalized
$T^1$-lifting gives comparisons of support and tangent space, with a
product description for surjective tangent localization
\cite[Theorems~3.8,~3.10]{Friedman}. Imagi studies concentration
of obstructions along curvilinear extensions
\cite[Theorem~1.2 and Definition~8.16]{Imagi25} and the singularity
invariants of generic fibres \cite[Theorems~1.7--1.8]{Imagi}. In the situation at hand the local
periods give equations for the actual image when the tangent image is a
proper subspace. Section~\ref{sec:examples} gives two compact examples with the same
reduced germ and tangent dimension. Their transverse rings are
$\C[x,y]/(x^2,xy,y^2)$ and $\C[x,y]/(x^2,y^2)$.
Their functors agree on curvilinear test rings: two square-zero
elements of $\C[t]/(t^m)$ have product zero by their orders in $t$.

The proof has two steps. Periods on the original links commute with
base change to Artin rings. Stokes' theorem then makes every linear
equation in \eqref{eq:kernel} hold on the full deformation space.
Conversely, a point of the proposed intersection over a small extension
supplies lifts of all the local germs. In Lemma~\ref{lem:obstruction}
we show that an actual global
obstruction which vanishes locally pairs trivially with a spanning set
in $H^1(X,\Omega^1_X)$. This gives a global lift, whose periods can be
adjusted by its Kodaira--Spencer class. The convergent map constructed
by the periods consequently identifies the full ideals.

\section{Periods on the original links}\label{sec:periods}

Let $B=\Def(X)$ and fix a representative of its proper flat family near
the given compact fibre. The morphism is Gorenstein after shrinking,
and its relative dualizing line commutes with base change. A generator
of $\omega_X$ extends to a generator $\Omega$ of this line, including
over the nonreduced base. Indeed $H^1(X,\omega_X)=0$ by duality and
$H^2(X,\OO_X)=0$. In a finite free complex computing the cohomology of
the relative dualizing line and its base changes, split the invertible
entries at the origin. The remaining complex has a rank-one term in
degree zero and zero terms in degrees $-1$ and $1$. Thus its generator
extends and properness permits us to remove the image of its zero
locus. Applied to $\OO$, the same argument shows that any other
generator differs by a unit on the base. Compare \cite[Lemma~3.1]{Friedman}.

On the punctured germ at $p$ a holomorphic relative three-form
determines a class on the oriented five-dimensional link. We use its
Poincar\'e dual in $V_p$. We give the details of its variation because
the possible Artin families include maps into a union of local
components which factor through none of them.

\begin{lemma}\label{lem:linkperiods}
Let $\mathcal U\to D_p$ be a marked semiuniversal family of one of the
local germs and let $D$ be a smooth parameter germ. Choose a relative
dualizing generator on $\mathcal U$ and multiply it by any holomorphic
unit on $\mathcal U\times D$. Its marked link class gives an analytic
isomorphism over $D$,
\[
       (D_p\times D,0)\longrightarrow (G_p\times D,0).
\]
It induces the cohomological link class on every Artin pullback. It
identifies the two components at a degree-six vertex and all the marked
local discriminant hyperplanes, with their ideals.
\end{lemma}
\begin{proof}
At an ordinary double point use the equation
$\sum_{j=1}^4 x_j^2=\tau$. The boundary three-cycle represents the
vanishing sphere in a smooth fibre. Its residue period is $\tau$ times
a unit. Scaling by $\sqrt{\tau}$ gives the spherical average of the
multiplier; odd terms integrate to zero and the constant term is
nonzero. This proves the assertion also in the presence of the extra
parameters.

At a degree-six vertex the analytic base has ring
$\C\{s,u,v\}/(su,sv)$ \cite[\S\S5--8]{Alt}. Work first on either
of its two smooth branches. The simultaneous modifications of
\cite[\S3, Proposition~3.8]{W5} have two marked nodes on the central
fibre over the line and three over the plane. Their total spaces are
respectively the total spaces of $\OO(-1,-1)^{\oplus2}$ over
$\PP^1\times\PP^1$ and $\OO(-1)^{\oplus3}$ over $\PP^2$. The map is
an isomorphism off the exceptional zero section, of codimension at
least two in each smooth total space. Since both maps to the branch
are flat Gorenstein, Hartogs extension for the pullback relative
generator and its inverse gives a relative generator on the incidence
space, also after taking its product with $D$. Let $q_j$ be the signed
nodal periods there. On the line both equation parameters have nonzero
derivative. On the plane the three are cyclic differences of three
diagonal entries modulo translation, so they have joint rank two at
the origin. Each $q_j$ is its equation parameter times a unit.

To identify the periods on the original link, use the cobordism between
that link and the links of the incidence nodes. On the central fibre
it is the region in the resolved neighborhood of $E_p$ obtained by
removing the inner small-resolution neighborhoods and truncating at
the outer boundary. It transports smoothly in the incidence family
by local triviality of its compact exterior with boundary. Each
$D_0\in H^2(E_p;\C)$ extends to the local resolution by retraction.
On this cobordism the boundary of the relative Poincar\'e dual of its
restriction has outer term $\operatorname{PD}(D_0|_{L_p})$ and inner
terms $(D_0\cdot C_j)\gamma_j$, where $C_j$ is the marked exceptional
curve at the inner node and $\gamma_j$ its signed link cycle. This
follows by restriction of Poincar\'e duality to each boundary component.
Stokes' theorem for the closed relative volume form gives the outer
period as $\sum_j(D_0\cdot C_j)q_j$. The circle-bundle sequence
identifies the outer class with $\sum_jq_j[C_j]$ in $H_2(E_p;\C)$.
Taking $D_0=K_{E_p}$ gives $\sum_j q_j=0$. The two or three marked
curves are the opposite pair or the alternating triple in the
introduction \cite[Proposition~3.7]{W5}. Their weighted sums identify
the sum-zero coefficient spaces with $V_{p,L}$ and $V_{p,P}$.

The resulting maps have invertible derivative along the respective
branches and vanish at their origins for all parameters in $D$.
Their inverses over $D$ therefore agree along $\{0\}\times D$.
The analytic ring of the union is the fibre product of the two branch
rings over $\OO_D$, so both maps and their inverses glue. A nonzero
plane parameter with one zero $q_j$ leaves a single ordinary double
point. The origin of either branch retains the original cone.

At a degree-five vertex the base is smooth of dimension four.
Write $E_p=\operatorname{Bl}_4\PP^2$ with line and exceptional classes
$H,e_1,\dots,e_4$. In the marking by the conic pencils put
\[
 f_i=H-e_i\ (1\leq i\leq4),\qquad f_5=2H-\textstyle\sum e_i,\qquad
 \lambda_i(c)=c\cdot f_i .
\]
These functionals identify $V_p=K_{E_p}^{\perp}$ with
$\{\sum\lambda_i=0\}\subset\C^5$. The local Grassmannian
family has five discriminant hyperplanes $\alpha_i(b)=0$ and one
residual node per vanishing $\alpha_i$ at a noncentral parameter
\cite[Proposition~3.1 and Lemma~3.2]{W5}. Its five marked boundary
cycles evaluate the link class by nonzero multiples of the $\lambda_i$.
For the homogeneous generator their periods are nonzero multiples of
$\alpha_i$ \cite[Lemma~3.4]{W5}. For a multiplier depending on $D$,
the period still vanishes on $\{\alpha_i=0\}\times D$ by the nodal
calculation on a dense open subset of that hyperplane. Holomorphic
division gives an $\alpha_i$ factor. Scaling the homogeneous generator
along a general line gives first derivative equal to the original
derivative multiplied by the value of the unit at the vertex. This is
the calculation of \cite[proof of Lemma~3.4]{W5}, with $D$ in place
of a disc. Consequently each remaining factor is a unit, the
derivative has rank four, and all five hyperplanes have the asserted
ideals.

We finish by checking the assertion over Artin rings. On the smooth
exterior of the central singularity in an Artin family over $A$, the
relative holomorphic de Rham complex resolves the constant sheaf $A$.
Indeed local relative polydisc coordinates exist by smoothness and
the radial Poincar\'e homotopy is coefficientwise. A relative closed
form thus defines a class in $H^3(L_p;\C)\otimes A$ on a small
punctured neighborhood, compatibly with restriction and base change.
For a smooth local base the analytic period just computed has this
class as its Taylor expansion. One way to see the compatibility is
to restrict the relative de Rham resolution of the inverse-image
sheaf of base functions to the compact boundary of a Milnor tube in
the smooth locus. That boundary is a differentiable product over a
small ball. Its degree-three direct image is the constant marked
cohomology tensored with the holomorphic base functions, computed
on a finite good cover of the link. Restriction of sheaves to the
boundary is exact, and in product charts the local relative de Rham
resolution commutes with evaluation in $A$. This proves the claim for
the Taylor jets, with any additional smooth parameters.

At the degree-six vertex apply this comparison to the Artin quotients
by powers of the maximal ideal of $\OO_{D_p\times D,0}$. Each such
quotient injects into the direct sum of the two corresponding branch
quotients: in the coordinates $(s,u,v)$ and the parameters of $D$ a
surviving monomial involves either $s$ or $u,v$, or neither. The
cohomological class and the Taylor jet of the glued analytic functions
agree in each branch quotient. Injectivity gives their equality on
the whole quotient. Every Artin point factors through one of these
quotients, proving the claim also on points outside either branch
functor.
\end{proof}

\begin{proposition}\label{prop:period}
The periods of $\Omega$ define an analytic morphism $P:B\to K\times_VG$.
Its differential has image $K$ and kernel $H^1(X,\Theta_X)$.
In each singular neighborhood it is the pullback of the local
isomorphism of Lemma~\ref{lem:linkperiods} over a smooth ambient germ
containing $B$.
\end{proposition}
\begin{proof}
Embed $B$ in a smooth germ $D$. Fix classifying maps $\phi_p:B\to
D_p$ with isomorphisms of the pulled-back local families. In the
family $\mathcal U_p\times D$ the restriction of the global family near $p$ is the closed
subgerm cut out by the graph of $\phi_p$. The ratio of the global
form to a pullback local generator is a holomorphic unit there.
It lifts to a holomorphic unit near the central point of
$\mathcal U_p\times D$. The preceding lemma gives a period morphism
over $D$ whose pullback is $P_p$. This construction respects the
fixed markings and records the dependence of the multiplier on global
parameters, even on a nonreduced base.

For every Artin point of $B$ the restricted global form gives one class
on the regular locus with coefficients in its base algebra. Its
restrictions to all the links are the classes just constructed.
Take the exterior in $Y$ of the local resolution neighborhoods.
For any degree-two cohomology class on $Y$, the Poincar\'e duals
of its restrictions to the boundaries together bound a relative
four-cycle in the exterior. Thus the corresponding sum of evaluations
of the global three-class is zero. With the oriented link identifications
this says precisely that the period tuple lies in the kernel
\eqref{eq:kernel}, with coefficients in the Artin algebra.
Apply this to every maximal-ideal jet of $B$. Separatedness of its
analytic local ring gives the linear equations of $K$ on $B$.
The equations of $G$ already hold on the individual local bases.

The derivative in directions tangent to $D$ at the vertex vanishes, as
all the central link classes are zero for every multiplier. The local
period derivatives on the local tangent spaces are isomorphisms.
Hence the full period map has precisely the kernel of tangent
localization. The local-to-global sequence and the marked
identification in \cite[Theorem~2.5]{W5} give that kernel and show
that the rank is $\dim K$. The linear equations just proved identify
its image with $K$ itself.
\end{proof}

\section{Lifting over small extensions}\label{sec:lifting}

The following obstruction calculation applies to isolated canonical
singularities beyond the local models of Theorem~\ref{thm:main}.

\begin{lemma}\label{lem:obstruction}
Let $X$ be a normal projective complex threefold with isolated Gorenstein
canonical singularities, $\omega_X\simeq\OO_X$ and $H^1(X,\OO_X)=0$.
Let $A'\twoheadrightarrow A$ be a small extension of local Artin
$\C$-algebras and let $X_A$ be a flat deformation of $X$.
If the restriction of $X_A$ at each singular point admits a lift to $A'$,
then $X_A$ admits a lift to $A'$.
\end{lemma}

\begin{proof}
Write $I$ for the kernel of the extension, so $\mathfrak m_{A'}I=0$.
We first explain the use of projective embeddings on the given Artin
deformation. Serre duality gives $H^2(X,\OO_X)=0$. The exact sequence of
units for a square-zero thickening therefore lifts any line bundle on $X$
to $X_A$, by induction on the length of $A$. For a very ample $M$ with
$H^1(X,M)=0$, a basis of $H^0(X,M)$ lifts to sections generating
the lifted bundle. They give a morphism $f_A:X_A\to\PP^N_A$ inducing
the embedding $f$ by $M$. This is a closed immersion: it is proper
with finite fibres, hence finite, and the map on structure sheaves is
surjective by Nakayama's lemma. In particular an analytic Artin
deformation can be treated algebraically, by GAGA for its ideal in
$\PP^N_A$. Local analytic lifts induce lifts of the completed local
rings. In comparing obstruction classes there we complete a
finite-type cotangent presentation, using continuous differentials.
Completion commutes with the finite free terms through each fixed
degree and is faithfully flat on the local obstruction module.

Put $L_X=L_{X/\C}$. The canonical obstruction to lifting $X_A$ is
\[
 o\in\Ext^2_X(L_X,\OO_X)\otimes I.
\]
Indeed the relative obstruction with coefficients in $I\otimes_\C\OO_X$ has
this target by flat base change for the cotangent complex and derived
adjunction. It is the composite of the transitivity boundary for
$X_A\to\Spec A\to\Spec A'$ with the shifted projection of
$L_{A/A'}$ onto $I[1]$. Its vanishing is equivalent to a flat lift
\cite[Tags 08UZ and 063Y]{Stacks}.
For the morphism $f_A$ above, functoriality of the transitivity triangle
identifies the image of $o$ under
\[
 \Ext^2_X(L_X,\OO_X)\otimes I
    \longrightarrow H^2(X,f^*T_{\PP^N})\otimes I
\]
with the pullback of the obstruction of $\PP^N_A$. Thus this image is
zero.

The negative cohomology sheaves of $L_X$ are supported on the finite
singular set $Z$. Finite-length modules over the local rings at these
points have grade three. The truncation of $L_X$ to $\Omega_X^1$ hence
induces an isomorphism in $\Ext^2(-,\OO_X)$, globally and locally
\cite[Lemma~2.2]{W5}. By Serre duality, $o$ is an $I$-valued functional
on $H^1(X,\Omega_X^1)$. The last vanishing says that this functional
kills the image of
\[
 H^1(X,f^*\Omega^1_{\PP^N})\longrightarrow H^1(X,\Omega_X^1).
\]
The connecting class of $1$ in the pulled-back Euler sequence maps to
$d\log M$, up to sign. Every line bundle is the difference, in
$\operatorname{Pic}(X)$, of two very ample bundles with the required
vanishing. Consequently $o$ annihilates $d\log\operatorname{Pic}(X)$.

By local duality, localization on $\Ext^2$ is dual to the map
\[
 H^1_Z(X,\Omega_X^1)\longrightarrow H^1(X,\Omega_X^1)
\]
\cite[Lemma~2.4(c)]{Gross}. At a singular point the image of $o$
with coefficients in the completed local ring is obtained by
precomposing the obstruction of the completed germ with the map
of cotangent complexes. The given local lift makes it zero, by the
same functoriality. Thus $o$ annihilates this second image. Write $j:X\setminus Z\hookrightarrow X$ and
$F=j_*\Omega^1_{X\setminus Z}$. The quotient by the displayed image is
$H^1(X,F)$. Indeed $H^1(X,F)$ embeds in $H^1(X\setminus Z,\Omega^1)$
by Leray, and $H^1(X,\Omega_X^1)\to H^1(X,F)$ is surjective since the
kernel and cokernel on sheaves have finite support. The local cohomology
sequence gives the assertion. Gross proves that the images of
$d\log\operatorname{Pic}(X)$ span $H^1(X,F)$
\cite[Lemma~2.5]{Gross}. Thus $o=0$.
\end{proof}

\begin{proof}[Proof of Theorem~\ref{thm:main}]
Write $B=\Def(X)$, let $G\subset V$ be the product of the marked local
period images, and put $T=\C^h\times(K\times_VG)$, with the indicated
scheme structure. Proposition~\ref{prop:period} gives the analytic
period map $P:B\to K\times_VG$. Its differential maps onto $K$ with
kernel $H^1(X,\Theta_X)$. Choose $h$ analytic functions $z$ vanishing
at the origin whose differentials restrict to a basis of the dual of this kernel. Then $F=(z,P):B\to T$ is a closed immersion of germs. Indeed
$dF$ is an isomorphism at the origin, and functions spanning the
cotangent space give a closed immersion by the analytic inverse
function theorem in a minimal smooth embedding.

Let $A'\twoheadrightarrow A$ be a small extension with kernel $I$.
Suppose $b\in B(A)$ and $y'\in T(A')$ restricts to $F(b)$. In the
construction of Proposition~\ref{prop:period}, lift the auxiliary
ambient parameter of $b$ to $A'$, which is possible since its ambient
germ is smooth. Over that lift, the inverse local period isomorphisms,
applied to the period coordinates of $y'$, give local hull parameters
extending those of $b$. In particular every singular germ of the
deformation over $b$ lifts to $A'$. Lemma~\ref{lem:obstruction} gives
a global lift, and the formal versality of $B$ gives an $A'$-point
$b'$ above $b$.

The difference of two lifts of the same point through a small
extension, viewed on local rings, is a derivation at the closed point
with values in $I$. Thus $y'-F(b')$ belongs to $T_0T\otimes I$.
Translate $b'$ by its inverse image under $dF\otimes 1_I$. The
resulting lift maps to $y'$. This proves formal smoothness of the
closed immersion $F$. More explicitly, successive lifting on the
Artin quotients of $\OO_{T,0}$ gives a section of $F$ on completions.
The completed quotient homomorphism is therefore an isomorphism.
Faithful flatness of completion gives the asserted analytic
isomorphism. The identifications on local branches and their marked
persistence loci follow from Proposition~\ref{prop:period}.
\end{proof}

\section{Which singularities can persist?}\label{sec:strata}

Choose one of the two local components at each degree-six vertex and
denote this choice by $S$. Let $V_S$ be the direct sum of the chosen
summands and the full spaces at all other points. Put $K_S=K\cap V_S$.
On $V_S$ take the following linear forms, with a separate index for
each marked cycle. At an original node take its coordinate. At a
degree-six vertex take the two or three signed nodal periods $q_j$ of
the chosen modification. At a degree-five vertex take the five forms
$\lambda_i$ of Lemma~\ref{lem:linkperiods}. Denote the resulting list
by $(\ell_j)_{j\in J_S}$. Its support at a point is the subset on
which the evaluations are nonzero.

\begin{corollary}\label{cor:strata}
The reduction of $\Def(X)$ is the product of $\C^h$ with the reduced
union of the $K_S$. The scheme-theoretic intersection of any nonempty collection
of selected bases equals $\C^h$ times the intersection of the corresponding
$K_S$. For $I\subset J_S$ put
\[
 F_I=K_S\cap\bigcap_{j\notin I}\ker\ell_j .
\]
The support $I$ occurs precisely when $\ell_i|_{F_I}\ne0$ for every
$i\in I$. Its stratum has dimension $h+\dim F_I$ and closure
$\C^h\times F_I$.
\end{corollary}
\begin{proof}
On taking the reduction of the intersection in Theorem~\ref{thm:main},
its support is the union of the linear sections $K_S$. The sum of the
ideals of any selected components is generated on $K$ by their defining
linear forms, since these already imply the local quadratic equations.
This gives the asserted intersection with its scheme structure.
Inside $F_I$ the required support is the complement of the kernels
of $\ell_i$ for $i\in I$. A complex vector space is never the union
of finitely many proper linear subspaces. This proves both existence and
density, including for empty $I$.
\end{proof}

Within a choice $S$ the closure contains exactly the feasible
subsets of $I$. Where two choices differ, their local link spaces meet
at zero, which represents the original cone. The preceding intersection
formula therefore also identifies strata represented on several choices.
A zero coordinate at an original node retains that node. At a nonzero
parameter for a vertex there is one residual node per vanishing marked
form, while the zero block retains the vertex itself.
Every feasible support is attained on a disc: take a line in $F_I$
through a point of that support and use Theorem~\ref{thm:main}.
The disc can be chosen projective over its base since an ample line
bundle on $X$ extends locally in the deformation when $H^2(X,\OO_X)=0$.

When only nodes and degree-six vertices occur, let $C_j$ be the
marked curves of the original and incidence nodes on the chosen
modification, represented on $Y$ as in Lemma~\ref{lem:linkperiods}.
In the coefficient coordinates $q_j$, $K_S$ is the kernel of
$(q_j)\mapsto\sum q_j[C_j]$ in $H_2(Y;\C)$. Indeed pairing with
each exceptional surface imposes the local sum-zero condition, and
then the remaining equation is \eqref{eq:kernel}. Thus, writing
$r_I=\dim\langle[C_i]:i\in I\rangle$, the dimension is
$h+|I|-r_I$. The existence criterion says that removing any single
$i\in I$ leaves $r_I$ unchanged. For a nodal threefold and $I$ the
set of all nodes, this is precisely the existence of a relation with
all coefficients nonzero in Friedman's criterion.

\section{Projective examples with different nilpotent structures}
\label{sec:examples}

We construct two partial resolutions of the same anticanonical
hypersurface. Their deformation spaces have smooth reductions of
dimension six and Zariski tangent spaces of dimension eight. The
nilradical has square zero in the first example and nonzero square in
the second.

Let $N=\{y\in\Z^5:\sum_i y_i=0\}$ and take the reflexive polytope
\[
  \Pi=[-1,1]^5\cap N_{\mathbb R}.
\]
Its thirty vertices are the permutations of $(1,1,-1,-1,0)$.
Fixing two coordinates with equal signs gives a unimodular triangle.
Fixing $y_i=1$, $y_j=-1$ gives a hexagon $H_{ij}$ whose remaining
coordinates are the permutations of $(1,0,-1)$. These are all the
two-dimensional faces. Each of the twenty hexagons has one interior
lattice point, its centre, and primitive edges. The polar has vertices
$\pm\bar e_i$ in $\Z^5/\Z(1,1,1,1,1)$, so the edge dual to each
hexagon is primitive.

Take a general anticanonical hypersurface in the toric fourfold
defined by the face fan. For a set $T$ of hexagons let $X_T$ be its
proper transform after star subdivision at the centres of all hexagons
outside $T$. These projective crepant subdivisions leave one exact
degree-six cone at each member of $T$. Indeed, the hypersurface meets
its one-dimensional orbit transversely in a single point, and the
transverse toric threefold is the cone of the marked hexagon. The unit
triangles and primitive edges give smooth germs elsewhere. Inserting
the remaining centres and taking a coherent maximal crepant
triangulation refining the resulting fan gives the required smooth
projective resolution $Y$. Its local restriction is the vertex blow-up
at each remaining cone \cite[\S8.1]{W4}.
The primal and polar lattice-point counts are $51$ and $11$,
respectively. Neither has a lattice point interior to a facet.
The dual-edge corrections in Batyrev's Hodge number formula vanish,
so $(h^{1,1}(Y),h^{2,1}(Y))=(46,6)$
\cite[Corollary~4.5.1]{Batyrev}.
The direct-image and higher-direct-image calculations for the tangent
sheaf at the vertex blow-up give $h=h^{2,1}(Y)=6$,
as in \cite[proof of Lemma~5.2]{W5}.

We take the ordered triples
\[
 A=(H_{54},H_{53},H_{51}),\qquad
 B=(H_{54},H_{24},H_{14}),\qquad
 C=(H_{24},H_{23},H_{21}).
\]
Thus $A$ and $B$ have the common face $H_{54}$, whereas $A$ and $C$
are disjoint. At each $H_{ij}$ read the free coordinates in increasing
order of their indices. Order its vertices $v_0,\ldots,v_5$ by the
successive triples
\[
 (-1,0,1),\ (0,-1,1),\ (1,-1,0),\ (1,0,-1),\ (0,1,-1),\ (-1,1,0).
\]
Write $D_k$ for the intersection vector of the $k$th boundary curve
against the ambient prime divisors. If $e_v$ denotes the coordinate of
the ray through $v$ and $c$ is the centre, the wall relation gives
\[
 D_k=e_{v_{k-1}}+e_{v_{k+1}}-e_{v_k}-e_c
 \qquad(k\bmod 6).
\]
Use the three rows $D_0-D_3,D_4-D_0,D_2-D_4$ for the coordinates
$(s,u,v)$. The first spans the line and the last two span the plane.
By \cite[Proposition~8.3]{W4} the ambient degrees detect all homology
relations here. This also follows from the vanishing correction term in
$h^{1,1}$, since the invariant divisors then span $H^2(Y;\Q)$.

\begin{proposition}\label{prop:cube-examples}
Put $R=\C\{w_1,\ldots,w_6\}$. The analytic deformation rings of
$X_{A\cup B}$ and $X_{A\cup C}$ are respectively
\[
 R\{x,y\}/(x^2,xy,y^2)
 \qquad\text{and}\qquad
 R\{x,y\}/(x^2,y^2).
\]
Every deformation over a reduced base is locally trivial at the
marked vertices in both examples.
\end{proposition}

\begin{proof}
The following table gives the unique relation, normalized by its first
entry, on each ordered triple. Entries are the $(s,u,v)$ coefficients
on successive faces. Extend each relation by zero on the other faces.
\[
\begin{array}{c|rrr|rrr|rrr}
 &\multicolumn{3}{c|}{1}&\multicolumn{3}{c|}{2}&\multicolumn{3}{c}{3}\\
 A &1&0&-1&-1&0&1&-1&-1&-1\\
 B &1&1&0&1&0&1&-1&0&-1\\
 C &1&-1&0&-1&1&0&1&-1&0
\end{array}
\]
The left kernels for $A\cup B$ and $A\cup C$ have the indicated
pairs of rows as bases. More explicitly, after dropping the centre
columns, which cancel, the five-face matrix has $15$ rows of rank
$13$, and the six-face matrix has $18$ rows of rank $16$. All entries
are given by the wall formula above.

On $A\cup B$ let $x,y$ be the coefficients of the two displayed
relations. The face $H_{51}$ gives $x^2=0$, and $H_{24}$ gives
$y^2=0$. The common face $H_{54}$ has block $(x+y,y,-x)$, giving
$xy=0$. On $A\cup C$, the two relations have disjoint support and
each has a nonzero product of a line and a plane coefficient on one
of its faces. The restricted quadrics thus span precisely $x^2,y^2$.
Theorem~\ref{thm:main} gives the two rings. In each case their
reductions are defined by $x=y=0$, where all local maps are constant.
\end{proof}

The nilradical in the second ring has cube zero and its square is
the free $R$-module generated by $xy$. Its successive quotients,
including the reduction, have ranks $(1,2,1)$ over $R$. The
corresponding ranks for the first ring are $(1,2)$. In particular, the distinction
between the two unrestricted germs persists after any analytic change
of coordinates.

\subsection{Two cones in the presence of original nodes}
For comparison consider the following partial resolutions, also obtained
by inserting selected hexagon centres. For the first, take in
$\Z\oplus\Z^3$ the convex hull of $(1,0)$, the six points
$(0,\pm e_i)$, and the sixteen points $(-1,a)$ where $a$ runs through
the permutations of $(\pm1,\pm1,0)$ and the sign triples with
$a_1a_2a_3=1$. Retain the hexagons with centres $(-1,e_1)$ and
$(-1,e_2)$. For the second take $X_+$ of \cite[\S5]{Partial},
with the two markings specified there. The two polytopes are entries
248 and 528 in the census accompanying \cite{W5}.

In both cases all primal edges and all dual edges of singular faces
are primitive. Every triangle is unimodular, every quadrilateral is
a unit square, and the remaining two-faces are degree-six hexagons.
The same coherent refinement as above resolves the retained vertices
and selects a small resolution at every square. The following counts
and ranks use the two retained hexagons in each case. Here $N_0$ is
the matrix of square circuits $e_i+e_j-e_k-e_l$, for opposite pairs
$(i,j),(k,l)$, and $M$ adjoins
the three rows for $(s,u,v)$ from each retained hexagon.
\[
\begin{array}{c|rrr|rr|rr}
 &\text{triangles}&\text{squares}&\text{hexagons}&
 h^{1,1}(Y)&h^{2,1}(Y)&\rk N_0&\rk M\\ \hline
 248&52&12&6&25&25&12&17\\
 528&32&30&4&27&18&20&22
\end{array}
\]
The accompanying computation reconstructs the faces from the ordered
vertices, checks unimodularity and the Hodge counts, and forms the
curve matrices from the wall formula and square circuits. The input
files contain only vertices and cyclic hexagon orders.

For the first threefold the unique relation has cone coefficients
$(1,1,1;-1,-1,-1)$. No choice of local branches contains this
direction. The reduced base is smooth of dimension $25$, and the
unrestricted ring is $\C\{w_1,\ldots,w_{25},t\}/(t^2)$ by
Theorem~\ref{thm:main}. In particular the nilpotent direction is
free over the entire reduced base.

For $X_+$ the kernel on the original nodes has dimension ten.
In the cyclic markings of \cite[\S5]{Partial}, the image on the two
cones has precisely the equations $u_0=0$ and $v_0+s_1=0$.
Writing $t=s_1$, $x=s_0$, $y=u_1$, $z=v_1$, the unrestricted ring is
\[
 \C\{w_1,\ldots,w_{28},t,x,y,z\}/(tx,ty,tz).
\]
The $28$ free coordinates include the ten original-node relations.
The two components have dimensions $31$ and $29$ and intersect along
the smooth germ of dimension $28$. A general point of the
$31$-dimensional component gives a smoothing. A general fibre on the other
component has one original node and one node near the first cone. In the ordering of \cite[\S5]{Partial}, the marked incidence relations
force precisely the period of original node $17$ and one of the three
plane periods to vanish on this component.

\begin{thebibliography}{Wue26c}
\bibitem[Alt97]{Alt} K.~Altmann, \emph{The versal deformation of an isolated toric Gorenstein singularity}, Invent. Math. \textbf{128} (1997), 443--479.
\bibitem[Bat94]{Batyrev} V.~V. Batyrev, \emph{Dual polyhedra and mirror symmetry for Calabi--Yau hypersurfaces in toric varieties}, J. Algebraic Geom. \textbf{3} (1994), 493--535.
\bibitem[Fri86]{Fri86} R.~Friedman, \emph{Simultaneous resolution of threefold double points}, Math. Ann. \textbf{274} (1986), 671--689.
\bibitem[Fri26]{Friedman} R.~Friedman, \emph{Unobstructed deformations for singular Calabi--Yau varieties}, \href{https://arxiv.org/abs/2506.09857v3}{arXiv:2506.09857v3} (2026).
\bibitem[Gro97]{Gross} M.~Gross, \emph{Deforming Calabi--Yau threefolds}, Math. Ann. \textbf{308} (1997), 187--220. Numbering follows the corrected version, \href{https://arxiv.org/abs/alg-geom/9506022v2}{arXiv:alg-geom/9506022v2} (2025).
\bibitem[Ima25]{Imagi25} Y.~Imagi, \emph{Deformations of compact Calabi--Yau conifolds}, \href{https://arxiv.org/abs/2504.03243v2}{arXiv:2504.03243v2} (2025).
\bibitem[Ima26]{Imagi} Y.~Imagi, \emph{Deformations of compact Calabi--Yau and Fano varieties with isolated singularities}, \href{https://arxiv.org/abs/2609.00758}{arXiv:2609.00758} (2026).
\bibitem[LLW18]{LLW} Y.-P.~Lee, H.-W.~Lin and C.-L.~Wang, \emph{Towards $A+B$ theory in conifold transitions for Calabi--Yau threefolds}, J. Differential Geom. \textbf{110} (2018), 495--541.
\bibitem[Nam02]{Nam} Y.~Namikawa, \emph{Stratified local moduli of Calabi--Yau threefolds}, Topology \textbf{41} (2002), 1219--1237.
\bibitem[Sta]{Stacks} The Stacks Project Authors, \emph{The Stacks Project}, Tags \href{https://stacks.math.columbia.edu/tag/08UZ}{08UZ} and \href{https://stacks.math.columbia.edu/tag/063Y}{063Y}.
\bibitem[Wue26a]{W4} B.~J. Wuebben, \emph{A vanishing-cycle obstruction to smoothing Calabi--Yau threefolds}, \href{https://arxiv.org/abs/2610.35911}{arXiv:2610.35911} (2026).
\bibitem[Wue26b]{W5} B.~J. Wuebben, \emph{Smoothing Calabi--Yau threefolds with nodes and del Pezzo cone points}, \href{https://arxiv.org/abs/2610.35915}{arXiv:2610.35915} (2026).
\bibitem[Wue26c]{Partial} B.~J. Wuebben, \emph{Partial smoothings of Calabi--Yau threefolds with degree-six del Pezzo cone points}, preprint (2026), \href{https://doi.org/10.5281/zenodo.24570171}{doi:10.5281/zenodo.24570171}.
\end{thebibliography}

\end{document}
